\documentclass[12pt]{article}
\usepackage[utf8]{inputenc}
\usepackage[T1]{fontenc}
\usepackage{hyperref}
\title{Fungwa Language Resources for K-12 Curriculum Development}
\author{Compiled by Claude Sonnet 5 (claude-sonnet-5) for Andr\'as Kornai}
\date{}
\begin{document}
\maketitle

\section{Introduction}
Fungwa (ISO 639-3: ula; also Cifungwa, Ura, Ula) is an endangered Kainji language spoken in Rafi Local Government Area, Niger State, Nigeria. Villages named in the sources include Sarari, Garbi, Jagaba and Ringa. Estimates of speaker numbers range from about 1,000 (Blench) to about 3,000. Speakers are reported to be shifting to Hausa.

\section{Primary Educational Resources}
No primers or school textbooks were found.

\section{Linguistic and Reference Materials}
\begin{itemize}
\item Wikipedia entry: \url{https://en.wikipedia.org/wiki/Fungwa_language}
\item Fungwa dictionary site on Webonary: \url{https://www.webonary.work/fungwa/?p=9444}
\item Samuel Akinbo (2019), work on vowel harmony, in Glottolog: \url{https://glottolog.org/resource/reference/id/568114}
\item Ifeoma Adiboshi (2018), work on tone and phonology, in Glottolog: \url{https://glottolog.org/resource/reference/id/568035}
\item Ifeoma Adiboshi (2014), in Glottolog: \url{https://glottolog.org/resource/reference/id/585678}
\item UBC Library open collections item (Fungwa-related academic work): \url{https://open.library.ubc.ca/collections/24/items/1.0401457}
\item Numerals page from the MPI Channumerals project: \url{https://lingweb.eva.mpg.de/channumerals/Fungwa.htm}
\end{itemize}
I only saw search snippets for these pages, so I cannot vouch for their full contents.

\section{Religious and Community Resources}
Joshua Project language profile: \url{https://joshuaproject.net/languages/ula}

\section{Digital and Audio Resources}
Documentation of Cifungwa Folktales (ELDP archive): \url{https://eldp.access.preservica.com/dk0436/}

\section{Curriculum Development Context}
The language has academic documentation (phonology, tone, vowel harmony, a folktale collection and what appears to be a dictionary site) but no found school materials. The folktale recordings and the dictionary could seed reading materials, subject to community consent. Pressure from Hausa is a risk to the language.
\end{document}